\section{A depth-blocking bound for general Boolean circuits}
\label{sec:boolean-blocking}

A simple circuit transformation slightly improves the exponential
dependence on growing intrinsic depth. It does not make the dependence
polynomial for general small circuits.

\begin{proposition}[Blocking Boolean layers]
\label{prop:boolean-layer-blocking}
Let $f$ be computed by an AND/OR/NOT circuit of fan-in at most two,
depth $h\ge1$, and at most $s$ gates and input wires. Put
\begin{align*}
 S&=\max\{16,2+d+s\},\qquad
 b=\max\{1,\lfloor\log_2\log_2 S\rfloor\},\\
 a&=\lfloor h/b\rfloor,\qquad r=h-ab,\qquad
 q=\max\{2,2a+\min\{2,r\}\}.
\end{align*}
Then $f$ has a depth-$q$ tanh representation with unit global margin,
width and physical weights bounded by a universal polynomial in $S$.
Consequently, for each fixed $\zeta>0$, at every learner depth $L\ge q$
the unrestricted posterior has a sample bound of the form
\[
 S^{\,O_\zeta(1+h/\log\log S)}
 \bigl(L+\varepsilon^{-1}+\log(2/\delta)\bigr)^{O_\zeta(1)},
\]
uniformly for every actual width $N\ge C_L(f)$ and every
$0<\sigma^2\le N^{-(1+\zeta)}$.
One has $q\le\max\{2,h\}$, so in particular the result applies
at every learner depth $L\ge\max\{2,h\}$.
For polynomial-size logarithmic-depth circuits, this is
$\exp[O_\zeta(\log^2 S/\log\log S)]$ times a polynomial in the
learner depth and statistical parameters.
\end{proposition}

\begin{proof}
First insert identity gates along skip wires and at an output of depth
less than $h$. The resulting layered circuit has polynomial size in $S$,
since $h\le s$ after removing unused gates. Divide its layers into
consecutive blocks of at most $b$ layers.

Each gate at the end of a block depends on at most $2^b$ signals at the
beginning of that block. Regard its subcomputation as a Boolean function
of these boundary signals, and express that function by its full DNF.
The DNF has at most
\[
 2^{2^b}\le S
\]
terms, each of length at most $2^b\le\log_2 S$. Implementing its
conjunctions in one threshold layer and its final disjunction in the
next uses polynomially many wires and integer coefficients bounded by
a polynomial in $S$. A gate whose function is constant can be
implemented by a constant threshold. Use two threshold layers for
every full block. A final one-layer block can simply retain its
original AND/OR/NOT threshold gates, rather than be expanded to two
layers; a longer final partial block uses two layers. The resulting
threshold depth is $2a+\min\{2,r\}\le h$. When $h=1$, pad to depth
two by inserting one identity threshold layer. All blocks together
therefore form a depth-$q$ threshold circuit of polynomial size and
polynomial integer weights. Choose odd integer thresholds after doubling literal
coefficients, so every gate has signed preactivation at least one on
Boolean inputs.

Here the threshold circuit already has small integer weights; no
threshold-to-majority replacement is needed. Convert its first $q-1$
layers directly to tanh. If $B\ge1$ bounds each gate's total absolute
incoming coefficients, including its bias, use gain
$\lambda=2\log(8B)$ and accuracy $1/(4B)$ at every hidden layer.
Inductively the affine perturbation from preceding approximate signs
is at most $1/4$, and the tanh error is at most
$2e^{-3\lambda/2}\le1/(4B)$. Biases are represented by an exact
hidden constant $1/2$, propagated using weight
$2\operatorname{arctanh}(1/2)$; its contribution is included in $B$.
At the last threshold, use twice the affine form as the linear readout.
It has signed value at least $3/2$. This construction has exactly $q$
weighted layers, no later biases, and polynomial width and weights.

The quantified fixed-intrinsic-depth theorem then gives a sample bound
$S^{O_\zeta(q)}$ times the displayed polynomial in $L$ and the
statistical parameters. Since $q=O(1+h/\log\log S)$, this is the
claimed estimate. Its small-width part is already expressed using the
actual complexity $C_L(f)$, so the alternative circuit's larger width
does not impose an additional lower-width assumption.
\end{proof}

This is a circuit-size/depth tradeoff rather than a polynomial-product
representation theorem. In particular, for balanced recursive majority
it improves the available quasipolynomial upper bound but does not
settle whether a polynomial bound is possible.
